\documentclass[11pt,a4paper]{article}

% ── Codificación, fuentes y lenguaje ───────────────────────────────────────────
\usepackage[utf8]{inputenc}
\usepackage[T1]{fontenc}
\usepackage[default,scale=0.95]{sourcesanspro}
\renewcommand{\familydefault}{\sfdefault}
\usepackage[spanish,es-noshorthands,es-tabla]{babel}

% ── Geometría y espaciado ─────────────────────────────────────────────────────
\usepackage[top=2.2cm, bottom=2.5cm, left=2.2cm, right=2.2cm, headheight=15pt]{geometry}
\usepackage{setspace}
\setstretch{1.18}
\usepackage{parskip}

% ── Matemáticas y unidades ────────────────────────────────────────────────────
\usepackage{amsmath}
\usepackage{amssymb}
\usepackage{siunitx}

% ── Circuitos ─────────────────────────────────────────────────────────────────
\usepackage[european, straightvoltages]{circuitikz}

% ── Tablas ────────────────────────────────────────────────────────────────────
\usepackage{booktabs}
\usepackage{tabularx}
\usepackage{array}
\usepackage{longtable}
\usepackage{colortbl}
\usepackage{enumitem}

% ── Colores, cajas e iconos ───────────────────────────────────────────────────
\usepackage[dvipsnames]{xcolor}
\usepackage{tcolorbox}
\tcbuselibrary{skins, breakable}
\usepackage{fontawesome5}
\usepackage{graphicx}

% ── Listados de código (dark-mode infográfico) ────────────────────────────────
%   Estilo base: fondo oscuro con números de línea. Los bloques lstlisting se
%   envuelven automáticamente en una caja tcolorbox redondeada.
\usepackage{listings}

% ── Secciones con barra de color (infografía) ────────────────────────────────
\usepackage{titlesec}
\titleformat{\section}
  {\normalfont\LARGE\bfseries\color{azulNoche}}
  {\colorbox{cyanNeon}{\color{fondoOscuro}\thesection}}{0.7em}{}
  [\vspace{2pt}\color{cyanNeon}\rule{\linewidth}{1.8pt}]
\titlespacing*{\section}{0pt}{24pt}{10pt}
\titleformat{\subsection}
  {\normalfont\large\bfseries\color{azulNoche}}
  {\textcolor{cyanNeon}{\thesubsection}}{0.7em}{}
  [\vspace{1pt}\color{grisLinea}\rule{\linewidth}{0.6pt}]
\titlespacing*{\subsection}{0pt}{14pt}{6pt}
\titleformat{\subsubsection}
  {\normalfont\normalsize\bfseries\color{azulNoche}}
  {\textcolor{cyanNeon}{\thesubsubsection}}{0.7em}{}
\titlespacing*{\subsubsection}{0pt}{10pt}{4pt}

% ── Cabeceras y pies de página ────────────────────────────────────────────────
\usepackage{fancyhdr}
\usepackage{lastpage}
\pagestyle{fancy}
\fancyhf{}
\renewcommand{\headrulewidth}{0pt}
\renewcommand{\footrulewidth}{0pt}
\fancyfoot[C]{%
  \begin{minipage}{\linewidth}
    \vspace{2pt}
    \color{cyanNeon}\rule{\linewidth}{1.2pt}\\[2pt]
    \centering\color{grisTexto}\footnotesize
    Página \thepage\ de \pageref{LastPage}
  \end{minipage}%
}

% ── Hiperenlaces ──────────────────────────────────────────────────────────────
\usepackage[colorlinks=true, linkcolor=azulNoche, urlcolor=cyanNeon]{hyperref}
\sloppy
\emergencystretch 3em

% ── Paleta de colores — tecnológico dark-mode ──────────────────────────────────
\definecolor{fondoOscuro}{HTML}{0D1B2A}     % Dark navy — fondo banda título
\definecolor{verdeTurquesa}{HTML}{004D40}    % Verde turquesa — bandas de marca
\definecolor{azulNoche}{HTML}{1B3A6B}       % Midnight blue — secciones, marcos
\definecolor{azulMedio}{HTML}{2A5298}       % Azul medio — variante de acento
\definecolor{cyanNeon}{HTML}{00C8E0}        % Cyan eléctrico — acentos primarios
\definecolor{verdeSignal}{HTML}{00B887}     % Verde señal — fortalezas, éxito
\definecolor{naranjaVivo}{HTML}{FF7043}     % Naranja vivo — mejoras, advertencias
\definecolor{rojoAlerta}{HTML}{E53935}      % Rojo alerta — errores
\definecolor{grisPapel}{HTML}{F4F6FA}       % Fondo tarjetas claras
\definecolor{grisLinea}{HTML}{C8D0E0}       % Bordes sutiles
\definecolor{grisTexto}{HTML}{6B7A99}       % Texto secundario
\definecolor{amarilloNota}{HTML}{FFB300}    % Notas / advertencias doradas

% Colores específicos para bloques de código (dark-mode)
\definecolor{codigoFondo}{HTML}{1E2D3D}      % Fondo del bloque de código
\definecolor{codigoComentario}{HTML}{7A8FA6} % Comentarios
\definecolor{codigoCadena}{HTML}{FF9D5C}     % Cadenas / strings
\definecolor{codigoKeyword}{HTML}{00C8E0}    % Palabras clave
\definecolor{codigoNumero}{HTML}{00E5A0}     % Números

% ── Banda de título: portada infográfica ──────────────────────────────────────
%   Diseño en 2 capas: banda de color (por defecto fondo oscuro) + franja de
%   acento cyan abajo. Uso: \bandaTitulo[color]{título}{subtítulo}.
%   El icono se puede cambiar con \renewcommand{\iconoBanda}{...} (FontAwesome).
\newcommand{\iconoBanda}{microchip}
\newcommand{\bandaTitulo}[3][fondoOscuro]{%
  \begin{tcolorbox}[
    enhanced,
    colback=#1, colframe=#1,
    boxrule=0pt, arc=8pt, outer arc=8pt,
    width=\linewidth,
    top=18pt, bottom=6pt, left=14pt, right=14pt,
    borderline south={3.5pt}{0pt}{cyanNeon},
  ]
    \begin{minipage}[c]{0.07\linewidth}
      \centering
      \textcolor{cyanNeon}{\fontsize{30}{30}\selectfont\faIcon{\iconoBanda}}
    \end{minipage}%
    \hspace{8pt}%
    \begin{minipage}[c]{0.88\linewidth}
      {\fontsize{20}{24}\selectfont\bfseries\color{white}#2}\par\vspace{5pt}%
      \textcolor{cyanNeon!80!white}{\small\faIcon{angle-right}~#3}%
    \end{minipage}
    \vspace{4pt}
  \end{tcolorbox}%
  \vspace{8pt}%
}

% ── Tarjetas de datos (infografía) ────────────────────────────────────────────
\tcbset{
  tarjetaDato/.style={
    enhanced, breakable,
    arc=6pt, outer arc=6pt,
    colback=grisPapel, colframe=azulNoche,
    boxrule=1pt,
    fonttitle=\bfseries\small, coltitle=white,
    attach boxed title to top left={yshift=-3mm, xshift=6mm},
    boxed title style={
      colback=azulNoche, colframe=azulNoche,
      arc=4pt, boxrule=0pt,
    },
    drop fuzzy shadow=azulNoche!30!white,
    top=8pt, bottom=8pt, left=10pt, right=10pt,
  },
  tarjetaIcono/.style={
    enhanced, breakable,
    arc=6pt, outer arc=6pt,
    colback=white, colframe=grisLinea, boxrule=0.8pt,
    fonttitle=\bfseries\small, coltitle=azulNoche,
    borderline west={3pt}{0pt}{cyanNeon},
    drop fuzzy shadow=azulNoche!25!white,
    top=6pt, bottom=6pt, left=10pt, right=10pt,
  },
  cajaTitulo/.style={
    enhanced, breakable,
    arc=6pt, outer arc=6pt,
    colback=azulNoche!10!grisPapel, colframe=azulNoche,
    boxrule=1pt,
    fonttitle=\bfseries, coltitle=white,
    attach boxed title to top left={yshift=-3mm, xshift=6mm},
    boxed title style={
      colback=azulNoche, colframe=azulNoche,
      arc=4pt, boxrule=0pt,
    },
    drop fuzzy shadow=azulNoche!25!white,
    top=8pt, bottom=8pt, left=10pt, right=10pt,
  },
  cajaContenido/.style={
    enhanced, breakable,
    arc=5pt, outer arc=5pt,
    colback=grisPapel, colframe=grisLinea,
    boxrule=0.8pt,
    fonttitle=\bfseries\small, coltitle=azulNoche,
    top=6pt, bottom=6pt, left=10pt, right=10pt,
  },
  cajaFortaleza/.style={
    enhanced, breakable,
    arc=6pt, outer arc=6pt,
    colback=verdeSignal!8!white, colframe=verdeSignal,
    boxrule=1pt,
    borderline west={4pt}{0pt}{verdeSignal},
    fonttitle=\bfseries\small, coltitle=verdeSignal!70!black,
    drop fuzzy shadow=verdeSignal!20!white,
    top=6pt, bottom=6pt, left=10pt, right=10pt,
  },
  cajaMejora/.style={
    enhanced, breakable,
    arc=6pt, outer arc=6pt,
    colback=naranjaVivo!8!white, colframe=naranjaVivo,
    boxrule=1pt,
    borderline west={4pt}{0pt}{naranjaVivo},
    fonttitle=\bfseries\small, coltitle=naranjaVivo!80!black,
    drop fuzzy shadow=naranjaVivo!20!white,
    top=6pt, bottom=6pt, left=10pt, right=10pt,
  },
  cajaRecomendacion/.style={
    enhanced, breakable,
    arc=6pt, outer arc=6pt,
    colback=cyanNeon!6!white, colframe=cyanNeon!60!azulNoche,
    boxrule=1pt,
    borderline west={4pt}{0pt}{cyanNeon},
    fonttitle=\bfseries\small, coltitle=azulNoche,
    drop fuzzy shadow=azulNoche!20!white,
    top=6pt, bottom=6pt, left=10pt, right=10pt,
  },
}

% ── Ícono + texto (encabezados de sección y elementos) ────────────────────────
\newcommand{\iconotexto}[2]{\textcolor{cyanNeon}{\faIcon{#1}}~#2}

% ── Listados de código: estilo dark + auto-envoltura ─────────────────────────
\lstdefinestyle{estiloCodigo}{
    backgroundcolor=\color{codigoFondo},
    basicstyle=\ttfamily\small\color{white},
    commentstyle=\color{codigoComentario}\itshape,
    stringstyle=\color{codigoCadena},
    keywordstyle=\color{codigoKeyword}\bfseries,
    numberstyle=\tiny\color{codigoComentario},
    numbers=left,
    numbersep=10pt,
    showstringspaces=false,
    breaklines=true,
    breakatwhitespace=false,
    tabsize=4,
    frame=none,
    captionpos=b,
    aboveskip=6pt,
    belowskip=4pt,
    xleftmargin=14pt,
    framexleftmargin=14pt,
    literate=
      {á}{{\'a}}1 {é}{{\'e}}1 {í}{{\'i}}1 {ó}{{\'o}}1 {ú}{{\'u}}1
      {Á}{{\'A}}1 {É}{{\'E}}1 {Í}{{\'I}}1 {Ó}{{\'O}}1 {Ú}{{\'U}}1
      {ñ}{{\~n}}1 {Ñ}{{\~N}}1 {ü}{{\"u}}1 {Ü}{{\"U}}1
      {¿}{{?`}}1 {¡}{{!`}}1
}
\lstdefinestyle{estiloPython}{
    style=estiloCodigo,
    language=Python,
}
\lstset{style=estiloCodigo}
\BeforeBeginEnvironment{lstlisting}{%
  \begin{tcolorbox}[
    enhanced, breakable,
    arc=6pt, outer arc=6pt,
    colback=codigoFondo, colframe=azulNoche,
    boxrule=1pt,
    drop fuzzy shadow=azulNoche!25!white,
    top=2pt, bottom=2pt, left=0pt, right=0pt,
  ]%
}
\AfterEndEnvironment{lstlisting}{\end{tcolorbox}}

% ── Cajas de contenido temáticas ──────────────────────────────────────────────
\newtcolorbox{cajaRecuerda}{
    enhanced, breakable,
    arc=6pt, outer arc=6pt,
    colback=azulNoche!10!grisPapel,
    colframe=azulNoche, boxrule=1pt,
    borderline west={4pt}{0pt}{cyanNeon},
    fonttitle=\bfseries\small\color{white},
    title={\faIcon{info-circle}~Recuerda},
    attach boxed title to top left={yshift=-3mm, xshift=6mm},
    boxed title style={colback=azulNoche, arc=4pt, boxrule=0pt},
    drop fuzzy shadow=azulNoche!25!white,
    left=10pt, right=10pt, top=8pt, bottom=8pt,
}

\newtcolorbox{cajaConcepto}{
    enhanced, breakable,
    arc=6pt, outer arc=6pt,
    colback=verdeSignal!8!white,
    colframe=verdeSignal, boxrule=1pt,
    borderline west={4pt}{0pt}{verdeSignal},
    fonttitle=\bfseries\small\color{white},
    title={\faIcon{lightbulb}~Concepto clave},
    attach boxed title to top left={yshift=-3mm, xshift=6mm},
    boxed title style={colback=verdeSignal!80!black, arc=4pt, boxrule=0pt},
    drop fuzzy shadow=verdeSignal!20!white,
    left=10pt, right=10pt, top=8pt, bottom=8pt,
}

\newtcolorbox{cajaEjemplo}{
    enhanced, breakable,
    arc=6pt, outer arc=6pt,
    colback=cyanNeon!5!white,
    colframe=cyanNeon!70!azulNoche, boxrule=1pt,
    borderline west={4pt}{0pt}{cyanNeon},
    fonttitle=\bfseries\small\color{fondoOscuro},
    title={\faIcon{code}~Ejemplo},
    attach boxed title to top left={yshift=-3mm, xshift=6mm},
    boxed title style={colback=cyanNeon, arc=4pt, boxrule=0pt},
    drop fuzzy shadow=azulNoche!20!white,
    left=10pt, right=10pt, top=8pt, bottom=8pt,
}

\newtcolorbox{cajaEjercicio}{
    enhanced, breakable,
    arc=6pt, outer arc=6pt,
    colback=azulMedio!7!white,
    colframe=azulMedio, boxrule=1pt,
    borderline west={4pt}{0pt}{azulMedio},
    fonttitle=\bfseries\small\color{white},
    title={\faIcon{pencil-alt}~Ejercicio},
    attach boxed title to top left={yshift=-3mm, xshift=6mm},
    boxed title style={colback=azulMedio, arc=4pt, boxrule=0pt},
    drop fuzzy shadow=azulNoche!20!white,
    left=10pt, right=10pt, top=8pt, bottom=8pt,
}

\newtcolorbox{cajaReto}{
    enhanced, breakable,
    arc=6pt, outer arc=6pt,
    colback=naranjaVivo!6!white,
    colframe=naranjaVivo, boxrule=1pt,
    borderline west={4pt}{0pt}{naranjaVivo},
    fonttitle=\bfseries\small\color{white},
    title={\faIcon{fire}~Reto},
    attach boxed title to top left={yshift=-3mm, xshift=6mm},
    boxed title style={colback=naranjaVivo, arc=4pt, boxrule=0pt},
    drop fuzzy shadow=naranjaVivo!20!white,
    left=10pt, right=10pt, top=8pt, bottom=8pt,
}

% ── Indicadores de dificultad ─────────────────────────────────────────────────
\newcommand{\facil}{\textcolor{verdeSignal}{\faIcon{circle}\,\textbf{Fácil}}}
\newcommand{\intermedio}{\textcolor{naranjaVivo}{\faIcon{circle}\,\textbf{Intermedio}}}
\newcommand{\dificil}{\textcolor{rojoAlerta}{\faIcon{circle}\,\textbf{Difícil}}}

% ─────────────────────────────────────────────────────────────────────────────

% ── Cabecera específica del reporte del skill ──
\renewcommand{\iconoBanda}{book}
\fancyhead[L]{\color{azulNoche}\small\bfseries \faIcon{book}~\textcolor{cyanNeon}{pallets\_itsdangerous}}
\fancyhead[R]{\color{grisTexto}\smallLibrería de firma y autenticación de datos}

\begin{document}
\bandaTitulo[fondoOscuro]{Skill: pallets\_itsdangerous}{"Usa este skill cuando necesites firmar criptográficamente datos, generar tokens seguros, serializar objetos con firma, implementar expiración de tokens, rotar claves secretas, o manejar datos en entornos no confiables. }

\section*{ItsDangerous - Firma y Autenticación de Datos}


\subsection*{Cuándo usar este skill}


Este skill debe usarse cuando se requiera:


\begin{itemize}[leftmargin=*,topsep=2pt,itemsep=1pt]
  \item Firmar datos para garantizar integridad y autenticidad en entornos no confiables (cookies, tokens de sesión, enlaces de verificación).
  \item Serializar objetos Python (JSON) y firmarlos para su transmisión segura.
  \item Implementar tokens con expiración temporal (enlaces de restablecimiento de contraseña, tokens de confirmación de email).
  \item Rotar claves secretas sin invalidar tokens existentes.
  \item Migrar algoritmos de firma manteniendo compatibilidad con tokens antiguos.
  \item Generar tokens seguros para URLs (Base64 URL-safe).
\end{itemize}


\subsection*{Estructura del proyecto}


El repositorio de ItsDangerous se organiza en los siguientes módulos principales:


\begin{itemize}[leftmargin=*,topsep=2pt,itemsep=1pt]
  \item \texttt{itsdangerous/\_\_init\_\_.py} - Punto de entrada, importa todas las clases públicas.
  \item \texttt{itsdangerous/signer.py} - Clases base \texttt{Signer} y \texttt{NoneAlgorithm}.
  \item \texttt{itsdangerous/serializer.py} - \texttt{Serializer} genérico.
  \item \texttt{itsdangerous/timed.py} - \texttt{TimestampSigner} y \texttt{TimedSerializer}.
  \item \texttt{itsdangerous/url\_safe.py} - \texttt{URLSafeSerializer} y \texttt{URLSafeTimedSerializer}.
  \item \texttt{itsdangerous/exc.py} - Jerarquía de excepciones.
  \item \texttt{itsdangerous/encoding.py} - Utilidades de codificación (base64, bytes).
\end{itemize}


\subsection*{API principal}


\subsubsection*{Firmas básicas}


\textbf{\texttt{Signer}} - Firma bytes y verifica la firma.


\begin{lstlisting}
class Signer:
    def __init__(
        self,
        secret_key: str | bytes | Iterable[str] | Iterable[bytes],
        salt: str | bytes | None = b"itsdangerous.Signer",
        sep: str | bytes = b".",
        key_derivation: str | None = None,
        digest_method: Any | None = None,
        algorithm: SigningAlgorithm | None = None,
    )
    def sign(self, value: str | bytes) -> bytes
    def unsign(self, signed_value: str | bytes) -> bytes  # lanza BadSignature
    def validate(self, signed_value: str | bytes) -> bool
\end{lstlisting}


\textbf{\texttt{SigningAlgorithm}} - Clase base para algoritmos de firma.


\begin{lstlisting}
class SigningAlgorithm:
    def get_signature(self, key: bytes, value: bytes) -> bytes
    def verify_signature(self, key: bytes, value: bytes, sig: bytes) -> bool
\end{lstlisting}


\textbf{\texttt{HMACAlgorithm}} - Algoritmo HMAC con hash configurable.


\begin{lstlisting}
class HMACAlgorithm(SigningAlgorithm):
    default_digest_method = staticmethod(_lazy_sha1)  # hashlib.sha1
    def __init__(self, digest_method=None)
\end{lstlisting}


\textbf{\texttt{NoneAlgorithm}} - Algoritmo de firma vacía (solo para pruebas).


\begin{lstlisting}
class NoneAlgorithm(SigningAlgorithm):
    def get_signature(self, key: bytes, value: bytes) -> bytes
    def verify_signature(self, key: bytes, value: bytes, sig: bytes) -> bool
\end{lstlisting}


\subsubsection*{Serialización}


\textbf{\texttt{Serializer}} - Serializa objetos (por defecto JSON) y los firma.


\begin{lstlisting}
class Serializer(Generic[_TSerialized]):
    default_serializer = json
    default_signer = Signer
    default_fallback_signers = []

    def __init__(
        self,
        secret_key: ...,
        salt: str | bytes | None = b"itsdangerous",
        serializer: ... = None,
        serializer_kwargs: dict | None = None,
        signer: type[Signer] | None = None,
        signer_kwargs: dict | None = None,
        fallback_signers: list | None = None,
    )
    def dumps(self, obj: Any, salt: ... = None) -> _TSerialized
    def loads(self, s: str | bytes, salt: ... = None, **kwargs) -> Any  # lanza BadSignature
    def loads_unsafe(self, s: str | bytes, salt: ... = None) -> tuple[bool, Any]
    def load_payload(self, payload: bytes, serializer: ... = None) -> Any
    def dump_payload(self, obj: Any) -> bytes
\end{lstlisting}


\subsubsection*{Firmas con timestamp}


\textbf{\texttt{TimestampSigner}} - Añade timestamp a la firma.


\begin{lstlisting}
class TimestampSigner(Signer):
    def get_timestamp() -> int
    def timestamp_to_datetime(ts: int) -> datetime
    def sign(self, value: str | bytes) -> bytes
    def unsign(self, signed_value: str | bytes, max_age: int | None = None, return_timestamp: bool = False) -> bytes | tuple[bytes, datetime]
    def validate(self, signed_value: str | bytes, max_age: int | None = None) -> bool
\end{lstlisting}


\textbf{\texttt{TimedSerializer}} - Serializa y añade timestamp.


\begin{lstlisting}
class TimedSerializer(Serializer):
    def loads(self, s: str | bytes, max_age: int | None = None, return_timestamp: bool = False, salt: ... = None) -> Any
    def loads_unsafe(self, s: str | bytes, max_age: int | None = None, salt: ... = None) -> tuple[bool, Any]
\end{lstlisting}


\subsubsection*{Serializadores URL-safe}


\textbf{\texttt{URLSafeSerializer}} - Serializa en formato seguro para URLs.


\begin{lstlisting}
class URLSafeSerializer(URLSafeSerializerMixin, Serializer[str]):
    def load_payload(self, payload: bytes, ...) -> Any
    def dump_payload(self, obj: Any) -> bytes
\end{lstlisting}


\textbf{\texttt{URLSafeTimedSerializer}} - Combina URL-safe con timestamp.


\begin{lstlisting}
class URLSafeTimedSerializer(URLSafeSerializerMixin, TimedSerializer[str]):
    pass
\end{lstlisting}


\subsubsection*{Utilidades de encoding}


\begin{lstlisting}
def want_bytes(s: str | bytes, encoding="utf-8", errors="strict") -> bytes
def base64_encode(string: str | bytes) -> bytes
def base64_decode(string: str | bytes) -> bytes
\end{lstlisting}


\subsubsection*{Excepciones}


\begin{lstlisting}
class BadData(Exception)           # base
class BadSignature(BadData)        # firma no coincide (payload disponible)
class BadTimeSignature(BadSignature)  # timestamp inválido (date_signed disponible)
class SignatureExpired(BadTimeSignature)  # expirado
class BadHeader(BadSignature)      # cabecera malformada
class BadPayload(BadData)          # payload no válido (original_error disponible)
\end{lstlisting}


\subsection*{Patrones de uso}


\subsubsection*{1. Firma básica con serializador URL-safe}


\begin{lstlisting}
from itsdangerous import URLSafeSerializer

s = URLSafeSerializer("secret-key", salt="auth")
token = s.dumps({"id": 5, "name": "itsdangerous"})
data = s.loads(token)
\end{lstlisting}


\subsubsection*{2. Tokens con expiración temporal}


\begin{lstlisting}
from itsdangerous import TimestampSigner

s = TimestampSigner("secret-key")
signed = s.sign("foo")
s.unsign(signed, max_age=5)  # lanza SignatureExpired si > 5 segundos
\end{lstlisting}


\subsubsection*{3. Rotación de claves secretas}


\begin{lstlisting}
SECRET_KEYS = ["old-key", "new-key"]
s = Serializer(SECRET_KEYS)
token = s.dumps({"id": 42})  # firma con "new-key"

# Rotar clave
SECRET_KEYS.append("newer-key")
del SECRET_KEYS[0]
s = Serializer(SECRET_KEYS)
s.loads(token)  # válido, se prueba con todas las claves
\end{lstlisting}


\subsubsection*{4. Migración de algoritmo con fallback signers}


\begin{lstlisting}
import hashlib

s = Serializer(
    "secret-key",
    signer_kwargs={"digest_method": hashlib.sha512},
    fallback_signers=[{"digest_method": hashlib.sha1}]
)
\end{lstlisting}


\subsubsection*{5. Manejo de errores con inspección de payload}


\begin{lstlisting}
from itsdangerous.exc import BadSignature, BadData

try:
    data = s.loads(token)
except BadSignature as e:
    if e.payload is not None:
        try:
            decoded = s.load_payload(e.payload)
        except BadData:
            pass  # payload corrupto
\end{lstlisting}


\subsubsection*{6. Carga insegura para debugging}


\begin{lstlisting}
sig_ok, payload = s.loads_unsafe(data)
\end{lstlisting}


\subsection*{Configuración}


\subsubsection*{Parámetros principales}


\begin{itemize}[leftmargin=*,topsep=2pt,itemsep=1pt]
  \item \textbf{\texttt{secret\_key}}: Clave secreta (bytes aleatorios). Puede ser lista para rotación.
  \item \textbf{\texttt{salt}}: Distingue contextos de firma. No necesita ser secreto.
  \item \textbf{\texttt{key\_derivation}}: Método para derivar clave (\texttt{"django-concat"}, \texttt{"concat"}, \texttt{"hmac"}, \texttt{"none"}).
  \item \textbf{\texttt{digest\_method}}: Función hash para HMAC (por defecto \texttt{hashlib.sha1}).
  \item \textbf{\texttt{serializer}}: Objeto con \texttt{dumps}/\texttt{loads} (por defecto \texttt{json}).
  \item \textbf{\texttt{fallback\_signers}}: Lista de configuraciones para verificar firmas antiguas.
  \item \textbf{\texttt{sep}}: Separador entre valor y firma (por defecto \texttt{"."}).
  \item \textbf{\texttt{max\_age}}: Edad máxima en segundos para tokens temporales.
  \item \textbf{\texttt{return\_timestamp}}: Si es \texttt{True}, \texttt{unsign}/\texttt{loads} retorna también el timestamp.
\end{itemize}


\subsection*{Ejemplos de código reutilizable}


\subsubsection*{Generar token de confirmación de email}


\begin{lstlisting}
from itsdangerous import URLSafeTimedSerializer

serializer = URLSafeTimedSerializer("secret-key", salt="email-confirm")
token = serializer.dumps({"user_id": 123})
# Enviar token por email...
\end{lstlisting}


\subsubsection*{Verificar token con expiración}


\begin{lstlisting}
from itsdangerous.exc import SignatureExpired, BadSignature

try:
    data = serializer.loads(token, max_age=3600)  # 1 hora
except SignatureExpired:
    print("El token ha expirado")
except BadSignature:
    print("Token inválido")
\end{lstlisting}


\subsubsection*{Rotar claves en aplicación Flask}


\begin{lstlisting}
from itsdangerous import URLSafeTimedSerializer

class TokenManager:
    def __init__(self, secret_keys):
        self.secret_keys = secret_keys
        self.serializer = URLSafeTimedSerializer(secret_keys)
    
    def rotate_key(self, new_key):
        self.secret_keys.insert(0, new_key)
        self.serializer = URLSafeTimedSerializer(self.secret_keys)
\end{lstlisting}


\subsection*{Glosario}


\begin{itemize}[leftmargin=*,topsep=2pt,itemsep=1pt]
  \item \textbf{Firma criptográfica}: Proceso de generar un código (HMAC) que verifica la integridad y autenticidad de los datos.
  \item \textbf{HMAC}: Hash-based Message Authentication Code, algoritmo que combina una clave secreta con un hash.
  \item \textbf{Salt}: Valor aleatorio que se añade a la clave para derivar claves diferentes en distintos contextos.
  \item \textbf{Key derivation}: Proceso de generar una clave criptográfica a partir de una clave maestra y un salt.
  \item \textbf{Timestamp}: Marca de tiempo Unix que indica cuándo se firmó el dato.
  \item \textbf{URL-safe}: Codificación Base64 que usa caracteres seguros para URLs (\texttt{-} y \texttt{\_} en lugar de \texttt{+} y \texttt{/}).
  \item \textbf{Fallback signers}: Configuraciones de firmantes antiguos que permiten verificar tokens firmados con algoritmos o claves anteriores.
  \item \textbf{Rotación de claves}: Práctica de cambiar periódicamente las claves secretas manteniendo las antiguas para validar tokens existentes.
\end{itemize}


\subsection*{Prerrequisitos}


\begin{itemize}[leftmargin=*,topsep=2pt,itemsep=1pt]
  \item Python 3.8 o superior.
  \item No tiene dependencias externas (solo stdlib).
  \item Instalación: \texttt{pip install itsdangerous}
\end{itemize}


\section*{api}

\section*{API de ItsDangerous}


\subsection*{Módulo \texttt{itsdangerous.signer}}


\subsubsection*{\texttt{Signer}}


Firma bytes y verifica la firma usando HMAC.


\begin{lstlisting}
class Signer:
    def __init__(
        self,
        secret_key: str | bytes | Iterable[str] | Iterable[bytes],
        salt: str | bytes | None = b"itsdangerous.Signer",
        sep: str | bytes = b".",
        key_derivation: str | None = None,
        digest_method: Any | None = None,
        algorithm: SigningAlgorithm | None = None,
    )
\end{lstlisting}


\textbf{Parámetros:}


\begin{itemize}[leftmargin=*,topsep=2pt,itemsep=1pt]
  \item \texttt{secret\_key}: Clave secreta. Puede ser una cadena, bytes, o iterable de claves para rotación.
  \item \texttt{salt}: Valor para derivar claves diferentes en distintos contextos. No necesita ser secreto.
  \item \texttt{sep}: Separador entre el valor y la firma (por defecto \texttt{"."}).
  \item \texttt{key\_derivation}: Método de derivación de clave (\texttt{"concat"}, \texttt{"django-concat"}, \texttt{"hmac"}, \texttt{"none"}).
  \item \texttt{digest\_method}: Función hash para HMAC (por defecto \texttt{hashlib.sha1}).
  \item \texttt{algorithm}: Instancia de \texttt{SigningAlgorithm} (por defecto \texttt{HMACAlgorithm}).
\end{itemize}


\textbf{Métodos:}


\begin{lstlisting}
def sign(self, value: str | bytes) -> bytes
\end{lstlisting}


Firma un valor y retorna \texttt{valor.separador.firma}.


\begin{lstlisting}
def unsign(self, signed_value: str | bytes) -> bytes
\end{lstlisting}


Verifica la firma y retorna el valor original. Lanza \texttt{BadSignature} si la firma es inválida.


\begin{lstlisting}
def validate(self, signed_value: str | bytes) -> bool
\end{lstlisting}


Retorna \texttt{True} si la firma es válida, \texttt{False} en caso contrario.


\textbf{Ejemplo:}


\begin{lstlisting}
from itsdangerous import Signer

s = Signer("secret-key")
signed = s.sign("hello")
# b'hello.XXXXX'
original = s.unsign(signed)
# b'hello'
\end{lstlisting}


\subsubsection*{\texttt{SigningAlgorithm}}


Clase base abstracta para algoritmos de firma.


\begin{lstlisting}
class SigningAlgorithm:
    def get_signature(self, key: bytes, value: bytes) -> bytes
    def verify_signature(self, key: bytes, value: bytes, sig: bytes) -> bool
\end{lstlisting}


\subsubsection*{\texttt{HMACAlgorithm}}


Implementación HMAC del algoritmo de firma.


\begin{lstlisting}
class HMACAlgorithm(SigningAlgorithm):
    default_digest_method = staticmethod(_lazy_sha1)
    
    def __init__(self, digest_method=None)
\end{lstlisting}


\textbf{Parámetros:}


\begin{itemize}[leftmargin=*,topsep=2pt,itemsep=1pt]
  \item \texttt{digest\_method}: Función hash (por defecto \texttt{hashlib.sha1}).
\end{itemize}


\subsubsection*{\texttt{NoneAlgorithm}}


Algoritmo de firma vacía que no genera firma real. Solo para pruebas.


\begin{lstlisting}
class NoneAlgorithm(SigningAlgorithm):
    def get_signature(self, key: bytes, value: bytes) -> bytes
    def verify_signature(self, key: bytes, value: bytes, sig: bytes) -> bool
\end{lstlisting}


\subsection*{Módulo \texttt{itsdangerous.serializer}}


\subsubsection*{\texttt{Serializer}}


Serializa objetos Python (por defecto JSON) y los firma.


\begin{lstlisting}
class Serializer(Generic[_TSerialized]):
    default_serializer = json
    default_signer = Signer
    default_fallback_signers = []

    def __init__(
        self,
        secret_key: str | bytes | Iterable[str] | Iterable[bytes],
        salt: str | bytes | None = b"itsdangerous",
        serializer: Any = None,
        serializer_kwargs: dict | None = None,
        signer: type[Signer] | None = None,
        signer_kwargs: dict | None = None,
        fallback_signers: list | None = None,
    )
\end{lstlisting}


\textbf{Parámetros:}


\begin{itemize}[leftmargin=*,topsep=2pt,itemsep=1pt]
  \item \texttt{secret\_key}: Clave secreta (igual que en \texttt{Signer}).
  \item \texttt{salt}: Salt para derivación de clave.
  \item \texttt{serializer}: Módulo u objeto con \texttt{dumps}/\texttt{loads} (por defecto \texttt{json}).
  \item \texttt{serializer\_kwargs}: Argumentos adicionales para el serializador.
  \item \texttt{signer}: Clase de firmante a usar (por defecto \texttt{Signer}).
  \item \texttt{signer\_kwargs}: Argumentos para el constructor del firmante.
  \item \texttt{fallback\_signers}: Lista de diccionarios con kwargs para firmantes de respaldo.
\end{itemize}


\textbf{Métodos:}


\begin{lstlisting}
def dumps(self, obj: Any, salt: str | bytes | None = None) -> _TSerialized
\end{lstlisting}


Serializa \texttt{obj} a JSON, lo firma y retorna el resultado.


\begin{lstlisting}
def loads(self, s: str | bytes, salt: str | bytes | None = None, **kwargs) -> Any
\end{lstlisting}


Verifica la firma y deserializa el payload. Lanza \texttt{BadSignature} si falla.


\begin{lstlisting}
def loads_unsafe(self, s: str | bytes, salt: str | bytes | None = None) -> tuple[bool, Any]
\end{lstlisting}


Retorna \texttt{(éxito, payload)} sin lanzar excepciones.


\begin{lstlisting}
def load_payload(self, payload: bytes, serializer: Any = None) -> Any
\end{lstlisting}


Deserializa un payload ya verificado.


\begin{lstlisting}
def dump_payload(self, obj: Any) -> bytes
\end{lstlisting}


Serializa un objeto a payload sin firmar.


\textbf{Ejemplo:}


\begin{lstlisting}
from itsdangerous import Serializer

s = Serializer("secret-key")
token = s.dumps({"user": "alice"})
# b'{"user":"alice"}.XXXXX'
data = s.loads(token)
# {"user": "alice"}
\end{lstlisting}


\subsection*{Módulo \texttt{itsdangerous.timed}}


\subsubsection*{\texttt{TimestampSigner}}


Firmante que añade timestamp a la firma.


\begin{lstlisting}
class TimestampSigner(Signer):
    def get_timestamp() -> int
    def timestamp_to_datetime(ts: int) -> datetime
    def sign(self, value: str | bytes) -> bytes
    def unsign(
        self,
        signed_value: str | bytes,
        max_age: int | None = None,
        return_timestamp: bool = False
    ) -> bytes | tuple[bytes, datetime]
    def validate(self, signed_value: str | bytes, max_age: int | None = None) -> bool
\end{lstlisting}


\textbf{Métodos:}


\begin{lstlisting}
def get_timestamp() -> int
\end{lstlisting}


Retorna la marca de tiempo Unix actual.


\begin{lstlisting}
def timestamp_to_datetime(ts: int) -> datetime
\end{lstlisting}


Convierte timestamp Unix a \texttt{datetime} UTC.


\begin{lstlisting}
def sign(self, value: str | bytes) -> bytes
\end{lstlisting}


Firma el valor con timestamp: \texttt{valor.sep.timestamp.sep.firma}.


\begin{lstlisting}
def unsign(
    self,
    signed_value: str | bytes,
    max_age: int | None = None,
    return_timestamp: bool = False
) -> bytes | tuple[bytes, datetime]
\end{lstlisting}


Verifica firma y timestamp. Si \texttt{max\_age} se supera, lanza \texttt{SignatureExpired}. Si \texttt{return\_timestamp=True}, retorna \texttt{(valor, datetime)}.


\begin{lstlisting}
def validate(self, signed_value: str | bytes, max_age: int | None = None) -> bool
\end{lstlisting}


Retorna \texttt{True} si la firma es válida y no ha expirado.


\textbf{Ejemplo:}


\begin{lstlisting}
from itsdangerous import TimestampSigner

s = TimestampSigner("secret-key")
signed = s.sign("foo")
# b'foo.1234567890.XXXXX'
original = s.unsign(signed, max_age=3600)
# b'foo'
\end{lstlisting}


\subsubsection*{\texttt{TimedSerializer}}


Serializador que añade timestamp.


\begin{lstlisting}
class TimedSerializer(Serializer):
    def loads(
        self,
        s: str | bytes,
        max_age: int | None = None,
        return_timestamp: bool = False,
        salt: str | bytes | None = None
    ) -> Any
    def loads_unsafe(
        self,
        s: str | bytes,
        max_age: int | None = None,
        salt: str | bytes | None = None
    ) -> tuple[bool, Any]
\end{lstlisting}


\textbf{Ejemplo:}


\begin{lstlisting}
from itsdangerous import TimedSerializer

s = TimedSerializer("secret-key")
token = s.dumps({"user": 42})
data = s.loads(token, max_age=300)
\end{lstlisting}


\subsection*{Módulo \texttt{itsdangerous.url\_safe}}


\subsubsection*{\texttt{URLSafeSerializerMixin}}


Mixin que añade codificación Base64 URL-safe y compresión zlib.


\begin{lstlisting}
class URLSafeSerializerMixin:
    def load_payload(self, payload: bytes, *args, **kwargs) -> Any
    def dump_payload(self, obj: Any) -> bytes
\end{lstlisting}


\subsubsection*{\texttt{URLSafeSerializer}}


Serializador con salida URL-safe.


\begin{lstlisting}
class URLSafeSerializer(URLSafeSerializerMixin, Serializer[str]):
    pass
\end{lstlisting}


\subsubsection*{\texttt{URLSafeTimedSerializer}}


Serializador URL-safe con timestamp.


\begin{lstlisting}
class URLSafeTimedSerializer(URLSafeSerializerMixin, TimedSerializer[str]):
    pass
\end{lstlisting}


\textbf{Ejemplo:}


\begin{lstlisting}
from itsdangerous import URLSafeTimedSerializer

s = URLSafeTimedSerializer("secret-key")
token = s.dumps({"session_id": "abc"})
# 'eyJzZXNzaW9uX2lkIjoiYWJjIn0.XXXXX.XXXXX'
data = s.loads(token, max_age=86400)
\end{lstlisting}


\subsection*{Módulo \texttt{itsdangerous.encoding}}


\subsubsection*{Funciones de utilidad}


\begin{lstlisting}
def want_bytes(s: str | bytes, encoding="utf-8", errors="strict") -> bytes
\end{lstlisting}


Convierte cadena a bytes si es necesario.


\begin{lstlisting}
def base64_encode(string: str | bytes) -> bytes
\end{lstlisting}


Codifica en Base64 URL-safe.


\begin{lstlisting}
def base64_decode(string: str | bytes) -> bytes
\end{lstlisting}


Decodifica Base64 URL-safe.


\subsection*{Módulo \texttt{itsdangerous.exc}}


\subsubsection*{Jerarquía de excepciones}


\begin{lstlisting}
class BadData(Exception):
    """Excepción base para datos inválidos."""

class BadPayload(BadData):
    """Payload no pudo ser deserializado."""
    def __init__(self, message, original_error=None)

class BadSignature(BadData):
    """Firma no coincide."""
    def __init__(self, message, payload=None, original_error=None)

class BadTimeSignature(BadSignature):
    """Timestamp inválido o malformado."""
    def __init__(self, message, payload=None, date_signed=None, original_error=None)

class SignatureExpired(BadTimeSignature):
    """Firma expirada por max_age."""
    def __init__(self, message, payload=None, date_signed=None, original_error=None)

class BadHeader(BadSignature):
    """Cabecera del payload malformada."""
\end{lstlisting}


\section*{configuracion}

\section*{Configuración de ItsDangerous}


\subsection*{Parámetros de configuración}


\subsubsection*{\texttt{Signer} y derivados}


\begin{center}
\begin{tabular}{@{}llll@{}}
\toprule
Parámetro & Tipo & Default & Descripción \\
\midrule
\texttt{secret\_key} & `str \textbackslash{} & bytes \textbackslash{} & Iterable[str] \textbackslash{} \\
\texttt{salt} & `str \textbackslash{} & bytes \textbackslash{} & None` \\
\texttt{sep} & `str \textbackslash{} & bytes` & \texttt{b"."} \\
\texttt{key\_derivation} & `str \textbackslash{} & None` & \texttt{None} (equivale a \texttt{"concat"}) \\
\texttt{digest\_method} & `callable \textbackslash{} & None` & \texttt{hashlib.sha1} \\
\texttt{algorithm} & `SigningAlgorithm \textbackslash{} & None` & \texttt{HMACAlgorithm()} \\
\bottomrule
\end{tabular}
\end{center}



\subsubsection*{\texttt{Serializer} y derivados}


\begin{center}
\begin{tabular}{@{}llll@{}}
\toprule
Parámetro & Tipo & Default & Descripción \\
\midrule
\texttt{secret\_key} & `str \textbackslash{} & bytes \textbackslash{} & Iterable[str] \textbackslash{} \\
\texttt{salt} & `str \textbackslash{} & bytes \textbackslash{} & None` \\
\texttt{serializer} & `module \textbackslash{} & object` & \texttt{json} \\
\texttt{serializer\_kwargs} & `dict \textbackslash{} & None` & \texttt{None} \\
\texttt{signer} & `type[Signer] \textbackslash{} & None` & \texttt{Signer} \\
\texttt{signer\_kwargs} & `dict \textbackslash{} & None` & \texttt{None} \\
\texttt{fallback\_signers} & `list \textbackslash{} & None` & \texttt{[]} \\
\bottomrule
\end{tabular}
\end{center}



\subsubsection*{\texttt{TimestampSigner} y \texttt{TimedSerializer}}


\begin{center}
\begin{tabular}{@{}llll@{}}
\toprule
Parámetro & Tipo & Default & Descripción \\
\midrule
\texttt{max\_age} & `int \textbackslash{} & None` & \texttt{None} \\
\texttt{return\_timestamp} & \texttt{bool} & \texttt{False} & Si \texttt{True}, \texttt{unsign()}/\texttt{loads()} retornan también el timestamp como \texttt{datetime} UTC. \\
\bottomrule
\end{tabular}
\end{center}



\subsubsection*{\texttt{URLSafeSerializer} y \texttt{URLSafeTimedSerializer}}


Estos serializadores aplican automáticamente:


\begin{itemize}[leftmargin=*,topsep=2pt,itemsep=1pt]
  \item Codificación Base64 URL-safe (caracteres \texttt{-} y \texttt{\_} en lugar de \texttt{+} y \texttt{/}).
  \item Compresión zlib si el payload comprimido es menor que el original (indicado con prefijo \texttt{b"."}).
\end{itemize}


No tienen parámetros de configuración adicionales.


\subsection*{Métodos de derivación de clave (\texttt{key\_derivation})}


\begin{itemize}[leftmargin=*,topsep=2pt,itemsep=1pt]
  \item \textbf{\texttt{"concat"}} (default implícito): Concatena \texttt{salt + b"signer" + secret\_key}.
  \item \textbf{\texttt{"django-concat"}}: Concatena \texttt{salt + b"signer"} y usa \texttt{secret\_key} como clave directa (compatible con Django).
  \item \textbf{\texttt{"hmac"}}: Deriva la clave usando HMAC con el salt como mensaje y \texttt{secret\_key} como clave.
  \item \textbf{\texttt{"none"}}: No deriva, usa \texttt{secret\_key} directamente.
\end{itemize}


\subsection*{Algoritmos de firma disponibles}


\begin{itemize}[leftmargin=*,topsep=2pt,itemsep=1pt]
  \item \textbf{\texttt{HMACAlgorithm}}: Algoritmo por defecto. Usa HMAC con la función hash especificada en \texttt{digest\_method}.
  \item \textbf{\texttt{NoneAlgorithm}}: No genera firma real. Solo para pruebas, nunca en producción.
\end{itemize}


\subsection*{Variables de entorno}


ItsDangerous no requiere variables de entorno. La configuración se realiza programáticamente.


\subsection*{Mejores prácticas de configuración}


\begin{enumerate}[leftmargin=*,topsep=2pt,itemsep=1pt]
  \item \textbf{Usar salts diferentes para cada contexto}: Evita que un token generado para un propósito sea válido en otro.
\end{enumerate}


\begin{lstlisting}
   session_s = URLSafeSerializer(secret, salt="session")
   email_s = URLSafeSerializer(secret, salt="email")
\end{lstlisting}


\begin{enumerate}[leftmargin=*,topsep=2pt,itemsep=1pt]
  \item \textbf{Rotar claves periódicamente}: Mantén una lista de claves, añadiendo nuevas al inicio y eliminando antiguas.
\end{enumerate}


\begin{lstlisting}
   SECRET_KEYS = ["current-key", "previous-key"]
\end{lstlisting}


\begin{enumerate}[leftmargin=*,topsep=2pt,itemsep=1pt]
  \item \textbf{Usar \texttt{fallback\_signers} para migrar algoritmos}: Permite cambiar el algoritmo de hash sin invalidar tokens existentes.
\end{enumerate}


\begin{lstlisting}
   Serializer(
       secret,
       signer_kwargs={"digest_method": hashlib.sha256},
       fallback_signers=[{"digest_method": hashlib.sha1}]
   )
\end{lstlisting}


\begin{enumerate}[leftmargin=*,topsep=2pt,itemsep=1pt]
  \item \textbf{Elegir el serializador adecuado}:
  \item \texttt{Serializer}: Para tokens que no van en URLs.
  \item \texttt{URLSafeSerializer}: Para tokens en URLs o cookies.
  \item \texttt{TimedSerializer} / \texttt{URLSafeTimedSerializer}: Cuando se necesita expiración.
\end{enumerate}


\begin{enumerate}[leftmargin=*,topsep=2pt,itemsep=1pt]
  \item \textbf{Configurar \texttt{max\_age} según el caso de uso}:
  \item Enlaces de restablecimiento de contraseña: 1 hora (3600).
  \item Tokens de sesión: 1 día (86400).
  \item Enlaces de confirmación de email: 7 días (604800).
\end{enumerate}


\section*{ejemplos}

\section*{Ejemplos de código de ItsDangerous}


\subsection*{Ejemplo 1: Firma básica con Signer}


\begin{lstlisting}
from itsdangerous import Signer

s = Signer("secret-key")
signed = s.sign(b"hello")
print(signed)  # b'hello.XXXXX'

original = s.unsign(signed)
print(original)  # b'hello'

# Validar sin excepción
if s.validate(signed):
    print("Firma válida")
\end{lstlisting}


\subsection*{Ejemplo 2: Serialización JSON con Serializer}


\begin{lstlisting}
from itsdangerous import Serializer

s = Serializer("secret-key")
token = s.dumps({"user": "alice", "id": 42})
print(token)  # b'{"user":"alice","id":42}.XXXXX'

data = s.loads(token)
print(data)  # {"user": "alice", "id": 42}
\end{lstlisting}


\subsection*{Ejemplo 3: Token con expiración usando TimestampSigner}


\begin{lstlisting}
from itsdangerous import TimestampSigner
from itsdangerous.exc import SignatureExpired

s = TimestampSigner("secret-key")
signed = s.sign("datos-sensibles")
print(signed)  # b'datos-sensibles.1234567890.XXXXX'

try:
    original = s.unsign(signed, max_age=5)  # 5 segundos
except SignatureExpired:
    print("Token expirado")
\end{lstlisting}


\subsection*{Ejemplo 4: Serializador URL-safe con tiempo}


\begin{lstlisting}
from itsdangerous import URLSafeTimedSerializer

s = URLSafeTimedSerializer("secret-key", salt="email-confirm")
token = s.dumps({"user_id": 123})
print(token)  # 'eyJ1c2VyX2lkIjoxMjN9.XXXXX.XXXXX'

data = s.loads(token, max_age=3600)  # 1 hora
print(data)  # {"user_id": 123}
\end{lstlisting}


\subsection*{Ejemplo 5: Rotación de claves}


\begin{lstlisting}
from itsdangerous import Serializer

# Claves iniciales
SECRET_KEYS = ["old-key", "new-key"]
s = Serializer(SECRET_KEYS)

# Firmar con "new-key"
token = s.dumps({"id": 42})

# Rotar clave
SECRET_KEYS.insert(0, "newer-key")
SECRET_KEYS.pop()  # eliminar "old-key"
s = Serializer(SECRET_KEYS)

# Verificar token antiguo
data = s.loads(token)  # funciona con "new-key"
print(data)  # {"id": 42}
\end{lstlisting}


\subsection*{Ejemplo 6: Migración de algoritmo con fallback}


\begin{lstlisting}
import hashlib
from itsdangerous import Serializer

s = Serializer(
    "secret-key",
    signer_kwargs={"digest_method": hashlib.sha512},
    fallback_signers=[{"digest_method": hashlib.sha1}]
)

# Firmar con SHA512
token = s.dumps({"data": "test"})

# Verificar (funciona con ambos algoritmos)
data = s.loads(token)
print(data)  # {"data": "test"}
\end{lstlisting}


\subsection*{Ejemplo 7: Manejo de errores con payload}


\begin{lstlisting}
from itsdangerous import Serializer
from itsdangerous.exc import BadSignature, BadPayload

s = Serializer("secret-key")
token = s.dumps({"user": "bob"})

# Simular token corrupto
try:
    data = s.loads(token + b"tampered")
except BadSignature as e:
    print(f"Firma inválida: {e}")
    if e.payload:
        try:
            decoded = s.load_payload(e.payload)
            print(f"Payload: {decoded}")
        except BadPayload:
            print("Payload corrupto")
\end{lstlisting}


\subsection*{Ejemplo 8: Carga insegura}


\begin{lstlisting}
from itsdangerous import Serializer

s = Serializer("secret-key")
token = s.dumps({"user": "alice"})

sig_ok, payload = s.loads_unsafe(token)
print(f"Válido: {sig_ok}, Payload: {payload}")
# Válido: True, Payload: {"user": "alice"}

# Con token inválido
sig_ok, payload = s.loads_unsafe("token-invalido")
print(f"Válido: {sig_ok}, Payload: {payload}")
# Válido: False, Payload: None
\end{lstlisting}


\subsection*{Ejemplo 9: Uso de salt para diferentes contextos}


\begin{lstlisting}
from itsdangerous import URLSafeSerializer

# Diferentes salts para diferentes propósitos
session_serializer = URLSafeSerializer("secret-key", salt="session")
email_serializer = URLSafeSerializer("secret-key", salt="email")

session_token = session_serializer.dumps({"user_id": 123})
email_token = email_serializer.dumps({"action": "verify"})

# No se puede usar un token en otro contexto
try:
    session_serializer.loads(email_token)  # Lanza BadSignature
except BadSignature:
    print("Token de email no válido para sesión")
\end{lstlisting}


\subsection*{Ejemplo 10: Obtener timestamp de un token}


\begin{lstlisting}
from itsdangerous import TimestampSigner

s = TimestampSigner("secret-key")
signed = s.sign("data")

# Obtener valor y timestamp
value, timestamp = s.unsign(signed, return_timestamp=True)
print(f"Valor: {value}")
print(f"Firmado: {timestamp}")  # datetime UTC
\end{lstlisting}


\section*{glosario}

\section*{Glosario de ItsDangerous}


\subsection*{Términos técnicos}


\subsubsection*{Firma criptográfica}


Proceso de generar un código de autenticación (MAC) que permite verificar que los datos no han sido alterados y provienen de una fuente confiable. ItsDangerous usa HMAC para este propósito.


\subsubsection*{HMAC (Hash-based Message Authentication Code)}


Algoritmo que combina una clave secreta con una función hash (como SHA1, SHA256) para generar un código de autenticación. Verifica tanto la integridad como la autenticidad de los datos.


\subsubsection*{Clave secreta (secret\_key)}


Valor criptográfico (bytes aleatorios) usado para firmar y verificar datos. Debe mantenerse en secreto. Puede ser una lista para soportar rotación de claves.


\subsubsection*{Salt}


Valor aleatorio o único que se añade a la clave secreta durante la derivación para generar claves diferentes en distintos contextos. No necesita ser secreto. Ejemplo: usar salt="session" para cookies de sesión y salt="email" para tokens de email.


\subsubsection*{Derivación de clave (key derivation)}


Proceso de generar una clave criptográfica a partir de una clave maestra y un salt. ItsDangerous soporta varios métodos: concat, django-concat, hmac, none.


\subsubsection*{Separador (sep)}


Carácter o bytes que separan el valor original de la firma en el resultado firmado. Por defecto es \texttt{"."}.


\subsubsection*{Payload}


Los datos originales (serializados) que se firman. En \texttt{Serializer}, es el resultado de \texttt{json.dumps()}. En caso de firma inválida, el payload puede estar disponible para inspección.


\subsubsection*{Timestamp}


Marca de tiempo Unix (segundos desde 1970-01-01 00:00:00 UTC) que indica cuándo se firmó el dato. Usado por \texttt{TimestampSigner} para validar expiración.


\subsubsection*{max\_age}


Edad máxima permitida en segundos para un token con timestamp. Si el token es más antiguo, se lanza \texttt{SignatureExpired}.


\subsubsection*{URL-safe}


Codificación que usa caracteres seguros para incluir en URLs sin necesidad de percent-encoding. ItsDangerous usa Base64 modificada: \texttt{-} en lugar de \texttt{+}, \texttt{\_} en lugar de \texttt{/}, y omite el padding \texttt{=}.


\subsubsection*{Base64}


Esquema de codificación que representa datos binarios en texto ASCII. ItsDangerous usa Base64 URL-safe para tokens que van en URLs.


\subsubsection*{Compresión zlib}


Algoritmo de compresión sin pérdida. Los serializadores URL-safe comprimen el payload con zlib si el resultado es menor, anteponiendo \texttt{b"."} para indicar compresión.


\subsubsection*{Fallback signers}


Lista de configuraciones de firmantes antiguos que permiten verificar tokens firmados con algoritmos o configuraciones anteriores. Facilita la migración de algoritmos sin invalidar tokens existentes.


\subsubsection*{Rotación de claves}


Práctica de seguridad que consiste en cambiar periódicamente las claves secretas. ItsDangerous soporta rotación manteniendo una lista de claves: la primera se usa para firmar, todas se usan para verificar.


\subsubsection*{BadSignature}


Excepción lanzada cuando la firma de un token no coincide con la esperada. Indica que los datos pueden haber sido manipulados o que la clave/salt es incorrecta.


\subsubsection*{BadTimeSignature}


Excepción lanzada cuando el timestamp de un token está ausente o malformado. Hereda de \texttt{BadSignature}.


\subsubsection*{SignatureExpired}


Excepción lanzada cuando un token con timestamp ha excedido el \texttt{max\_age} especificado. Hereda de \texttt{BadTimeSignature}.


\subsubsection*{BadPayload}


Excepción lanzada cuando el payload no puede ser deserializado (ej. JSON inválido).


\subsubsection*{loads\_unsafe}


Método que retorna una tupla \texttt{(éxito, payload)} sin lanzar excepciones. Útil para debugging o cuando se necesita inspeccionar el contenido sin manejar errores.


\subsubsection*{Serializador}


Componente que convierte objetos Python a un formato de bytes (ej. JSON) y viceversa. En ItsDangerous, el serializador por defecto es \texttt{json}.


\subsubsection*{Signer}


Componente que firma bytes y verifica firmas. Es la clase base para todos los firmantes en ItsDangerous.


\subsubsection*{Algoritmo de firma (SigningAlgorithm)}


Clase abstracta que define cómo se genera y verifica una firma. Implementaciones: \texttt{HMACAlgorithm} (por defecto) y \texttt{NoneAlgorithm} (pruebas).


\section*{patrones}

\section*{Patrones de uso de ItsDangerous}


\subsection*{1. Firma y verificación básica}


\textbf{Problema:} Necesitas enviar datos a un cliente (ej. cookie) y verificar que no han sido manipulados al recibirlos de vuelta.


\textbf{Solución:} Usar \texttt{Signer} para firmar los datos.


\textbf{Flujo:}


\begin{enumerate}[leftmargin=*,topsep=2pt,itemsep=1pt]
  \item Crear instancia de \texttt{Signer} con clave secreta y salt único para el contexto.
  \item Firmar los datos con \texttt{sign()}.
  \item Enviar datos firmados al cliente.
  \item Al recibir, verificar con \texttt{unsign()} o \texttt{validate()}.
  \item Si \texttt{unsign()} lanza \texttt{BadSignature}, los datos fueron manipulados.
\end{enumerate}


\begin{lstlisting}
from itsdangerous import Signer

# 1. Configurar firmante
signer = Signer("clave-secreta", salt="cookie-session")

# 2. Firmar datos
signed_data = signer.sign("user_id=123")

# 3. Enviar al cliente (ej. en cookie)

# 4. Verificar al recibir
try:
    original = signer.unsign(signed_data)
    print(f"Datos válidos: {original}")
except BadSignature:
    print("Datos manipulados")
\end{lstlisting}


\subsection*{2. Tokens con expiración temporal}


\textbf{Problema:} Generar tokens que expiran después de un tiempo (enlaces de restablecimiento de contraseña, confirmación de email).\textbackslash{}m \textbf{Solución:} Usar \texttt{TimestampSigner} o \texttt{TimedSerializer} con \texttt{max\_age}.


\textbf{Flujo:}\textbackslash{}1. Crear \texttt{TimestampSigner} o \texttt{TimedSerializer}.


\begin{enumerate}[leftmargin=*,topsep=2pt,itemsep=1pt]
  \item Firmar/serializar datos con \texttt{sign()} o \texttt{dumps()}.
  \item Enviar token al usuario.
  \item Al recibir, verificar con \texttt{unsign()} o \texttt{loads()} especificando \texttt{max\_age}.
  \item Capturar \texttt{SignatureExpired} para tokens expirados.
\end{enumerate}


\begin{lstlisting}
from itsdangerous import TimestampSigner
from itsdangerous.exc import SignatureExpired, BadTimeSignature

# 1. Configurar
signer = TimestampSigner("clave-secreta", salt="reset-password")

# 2. Generar token
token = signer.sign("user_id=123")

# 3. Enviar por email...

# 4. Verificar con expiración de 1 hora
try:
    user_id = signer.unsign(token, max_age=3600)
    print(f"Token válido: {user_id}")
except SignatureExpired:
    print("Token expirado, solicitar nuevo")
except BadTimeSignature:
    print("Token inválido")
\end{lstlisting}


\subsection*{3. Serialización de objetos con firma}


\textbf{Problema:} Necesitas enviar estructuras de datos complejas (dicts, listas) de forma segura.


\textbf{Solución:} Usar \texttt{Serializer} con JSON.


\textbf{Flujo:}


\begin{enumerate}[leftmargin=*,topsep=2pt,itemsep=1pt]
  \item Crear \texttt{Serializer} con clave secreta.
  \item Serializar objeto con \texttt{dumps()}.
  \item Transmitir el token.
  \item Deserializar con \texttt{loads()}.
  \item Manejar \texttt{BadSignature} para datos inválidos.
\end{enumerate}


\begin{lstlisting}
from itsdangerous import Serializer

# 1. Configurar
serializer = Serializer("clave-secreta", salt="session")

# 2. Serializar y firmar
token = serializer.dumps({"user_id": 123, "role": "admin"})

# 3. Enviar token...

# 4. Deserializar
try:
    data = serializer.loads(token)
    print(f"Datos: {data}")
except BadSignature:
    print("Token inválido")
\end{lstlisting}


\subsection*{4. Rotación de claves secretas}


\textbf{Problema:} Necesitas cambiar la clave secreta periódicamente sin invalidar tokens existentes.


\textbf{Solución:} Pasar una lista de claves a \texttt{Signer}/\texttt{Serializer}. La primera clave se usa para firmar, todas se prueban para verificar.


\textbf{Flujo:}


\begin{enumerate}[leftmargin=*,topsep=2pt,itemsep=1pt]
  \item Inicializar con lista de claves (nueva primero).
  \item Firmar con la primera clave.
  \item Al rotar, añadir nueva clave al inicio y eliminar la más antigua.
  \item Tokens antiguos se verifican con claves anteriores.
\end{enumerate}


\begin{lstlisting}
from itsdangerous import Serializer

# 1. Claves iniciales
SECRET_KEYS = ["old-key", "new-key"]
serializer = Serializer(SECRET_KEYS)

# 2. Firmar con "new-key"
token = serializer.dumps({"id": 42})

# 3. Rotar clave
SECRET_KEYS.insert(0, "newer-key")
SECRET_KEYS.pop()  # eliminar "old-key"
serializer = Serializer(SECRET_KEYS)

# 4. Verificar token antiguo (funciona con "new-key")
data = serializer.loads(token)  # válido
\end{lstlisting}


\subsection*{5. Migración de algoritmo de firma}


\textbf{Problema:} Cambiar el algoritmo de hash (ej. SHA1 a SHA512) manteniendo compatibilidad con tokens antiguos.


\textbf{Solución:} Usar \texttt{fallback\_signers} para definir configuraciones de firmantes antiguos.


\textbf{Flujo:}


\begin{enumerate}[leftmargin=*,topsep=2pt,itemsep=1pt]
  \item Configurar \texttt{Serializer} con nuevo algoritmo en \texttt{signer\_kwargs}.
  \item Pasar \texttt{fallback\_signers} con configuraciones antiguas.
  \item Tokens nuevos se firman con algoritmo nuevo.
  \item Tokens antiguos se verifican con fallback signers.
\end{enumerate}


\begin{lstlisting}
import hashlib
from itsdangerous import Serializer

# 1. Configurar con nuevo algoritmo y fallback
serializer = Serializer(
    "secret-key",
    signer_kwargs={"digest_method": hashlib.sha512},
    fallback_signers=[{"digest_method": hashlib.sha1}]
)

# 2. Firmar nuevo token (SHA512)
token = serializer.dumps({"data": "test"})

# 3. Token antiguo (SHA1) también se verifica
old_token = "..."  # firmado con SHA1
data = serializer.loads(old_token)  # válido por fallback
\end{lstlisting}


\subsection*{6. Tokens URL-safe para enlaces}


\textbf{Problema:} Generar tokens que puedan incluirse en URLs sin caracteres problemáticos.


\textbf{Solución:} Usar \texttt{URLSafeSerializer} o \texttt{URLSafeTimedSerializer}.


\textbf{Flujo:}


\begin{enumerate}[leftmargin=*,topsep=2pt,itemsep=1pt]
  \item Crear serializador URL-safe.
  \item Serializar y firmar.
  \item Incluir token en URL.
  \item Al recibir, deserializar.
\end{enumerate}


\begin{lstlisting}
from itsdangerous import URLSafeTimedSerializer

# 1. Configurar
serializer = URLSafeTimedSerializer("secret-key", salt="email-verify")

# 2. Generar token
token = serializer.dumps({"user_id": 123})
# 'eyJ1c2VyX2lkIjoxMjN9.XXXXX.XXXXX'

# 3. Construir URL
url = f"https://example.com/verify?token={token}"

# 4. Verificar al recibir
try:
    data = serializer.loads(token, max_age=86400)
    print(f"Usuario verificado: {data['user_id']}")
except SignatureExpired:
    print("Enlace expirado")
\end{lstlisting}


\subsection*{7. Manejo robusto de errores con inspección de payload}


\textbf{Problema:} Necesitas diagnosticar por qué falló la verificación (firma inválida vs payload corrupto).


\textbf{Solución:} Capturar excepciones específicas y usar \texttt{payload} para inspeccionar.


\textbf{Flujo:}


\begin{enumerate}[leftmargin=*,topsep=2pt,itemsep=1pt]
  \item Intentar \texttt{loads()}.
  \item Capturar \texttt{BadSignature} y verificar \texttt{e.payload}.
  \item Si hay payload, intentar deserializarlo con \texttt{load\_payload()}.
  \item Capturar \texttt{BadPayload} si el payload está corrupto.
\end{enumerate}


\begin{lstlisting}
from itsdangerous import Serializer
from itsdangerous.exc import BadSignature, BadPayload

serializer = Serializer("secret-key")

try:
    data = serializer.loads(token)
except BadSignature as e:
    if e.payload is not None:
        try:
            # Intentar recuperar payload aunque firma sea inválida
            decoded = serializer.load_payload(e.payload)
            print(f"Payload recuperado: {decoded}")
        except BadPayload:
            print("Payload corrupto")
    else:
        print("Token sin payload")
\end{lstlisting}


\subsection*{8. Carga insegura para debugging}


\textbf{Problema:} Necesitas inspeccionar el contenido de un token sin que lance excepciones.


\textbf{Solución:} Usar \texttt{loads\_unsafe()} que retorna \texttt{(éxito, payload)}.


\begin{lstlisting}
from itsdangerous import Serializer

serializer = Serializer("secret-key")

sig_ok, payload = serializer.loads_unsafe(token)
if sig_ok:
    print(f"Token válido: {payload}")
else:
    print(f"Token inválido, payload: {payload}")
\end{lstlisting}


\section*{prerrequisitos}

\section*{Prerrequisitos de ItsDangerous}


\subsection*{Requisitos del sistema}


\begin{itemize}[leftmargin=*,topsep=2pt,itemsep=1pt]
  \item \textbf{Python}: 3.8 o superior.
  \item \textbf{Sistema operativo}: Cualquier sistema que soporte Python (Windows, macOS, Linux).
\end{itemize}


\subsection*{Dependencias}


ItsDangerous \textbf{no tiene dependencias externas}. Solo utiliza módulos de la biblioteca estándar de Python:


\begin{itemize}[leftmargin=*,topsep=2pt,itemsep=1pt]
  \item \texttt{hashlib} - Para funciones hash (SHA1, SHA256, SHA512, etc.).
  \item \texttt{hmac} - Para HMAC (Hash-based Message Authentication Code).
  \item \texttt{json} - Para serialización JSON (serializador por defecto).
  \item \texttt{base64} - Para codificación Base64 URL-safe.
  \item \texttt{zlib} - Para compresión de payloads en serializadores URL-safe.
  \item \texttt{datetime} - Para manejo de timestamps.
  \item \texttt{time} - Para obtener timestamps Unix.
  \item \texttt{typing} - Para anotaciones de tipo.
\end{itemize}


\subsection*{Instalación}


\subsubsection*{Desde PyPI (recomendado)}


\begin{lstlisting}
pip install itsdangerous
\end{lstlisting}


\subsubsection*{Desde el repositorio}


\begin{lstlisting}
git clone https://github.com/pallets/itsdangerous.git
cd itsdangerous
pip install -e .
\end{lstlisting}


\subsubsection*{Verificar instalación}


\begin{lstlisting}
import itsdangerous
print(itsdangerous.__version__)
\end{lstlisting}


\subsection*{Compatibilidad}


\begin{itemize}[leftmargin=*,topsep=2pt,itemsep=1pt]
  \item Compatible con Python 3.8, 3.9, 3.10, 3.11, 3.12.
  \item No tiene dependencias de sistema operativo.
  \item Puede usarse en cualquier framework web (Flask, Django, FastAPI, etc.) o aplicación Python.
\end{itemize}


\subsection*{Testing}


Para ejecutar los tests del proyecto:


\begin{lstlisting}
pip install -e ".[dev]"
pytest
\end{lstlisting}


\subsection*{Licencia}


BSD-3-Clause License. Ver archivo LICENSE en el repositorio.


\end{document}
